Die Arbeit ist im Kontext der bestehenden SWR-LS-Pipeline angesiedelt, die zur automatisierten Erstellung von Texten in LS+ eingesetzt wird. 
Diese Pipeline umfasst unter anderem Module zur Faktenextraktion, zum Pre-Editing, zur Generierung von LS+ Texten sowie zur Validierung der generierten Inhalte.
Das in dieser Arbeit entwickelte Evaluations- und Verbesserungssystem ist als eigenständiges Modul konzipiert, das nach der initialen Textgenerierung in die bestehende Pipeline integriert werden kann. 
Ziel ist es, den generierten Text nach der Erstellung automatisiert zu analysieren, anhand definierter Metriken zu bewerten und bei Bedarf iterativ zu verbessern.
Die Erweiterung soll als Qualitätsabsicherung dienen um die Verlässlichkeit der Prozesse zu erhöhen, ohne dabei grundlegende Strukturen zu verändern.